\documentclass[11pt,a4paper]{article}
\usepackage{amsmath,amssymb,amsthm}
\usepackage[margin=2.5cm]{geometry}
\usepackage{hyperref}

\newtheorem{definition}{Definition}[section]
\newtheorem{theorem}[definition]{Theorem}
\newtheorem{proposition}[definition]{Proposition}
\newtheorem{lemma}[definition]{Lemma}
\newtheorem{corollary}[definition]{Corollary}
\newtheorem{conjecture}[definition]{Conjecture}
\newtheorem{remark}[definition]{Remark}

\title{Hyperseed Formalization:\\
  Let's Get Chemical}
\author{ProtomegaTron (automated formalization)}
\date{2026-09-18}

\begin{document}
\maketitle

\begin{abstract}
We formalize the intellectual structure of Ben Goertzel's Substack article
``Let's Get Chemical'' (2026-05-05) within the Hyperseed ontology. The
article reports computational experiments in algorithmic chemistry ---
Lane-style protocell models with RAF sets, membrane conversion, and
evolutionary dynamics. The central contribution is ``protected release''
(value-relative anti-precedence): penalize excess canalization, not
repetition itself. Key mappings: protected release as fiber-selective
eviction, identity vs.\ function as section vs.\ fiber, conversion
bridge as transition function, value-justified share as fiber-weighted
allocation, and biphasic memory failure as indiscriminate fiber eviction.
\end{abstract}

\tableofcontents

%% =================================================================
\section{Protected release as fiber-selective eviction}
\label{sec:release}
%% =================================================================

\begin{theorem}[Value-relative anti-precedence --- claims 7--12, 21]
\label{thm:release}
The article's central conceptual contribution: don't penalize repetition
broadly (that kills individuation). Penalize \emph{excess canalization}:
when a pathway's actual share exceeds its value-justified share.

The rule:
\begin{enumerate}
\item For each pathway family $g$, compute its actual share $S_g(t)$ of
  the recent ecology.
\item Compute a value-justified share $S^*_g(t)$ via softmax over
  per-family value scores (membrane-conversion payoff per unit cost).
\item Define excess canalization $O_g = \max(S_g - S^*_g - \theta, 0)$.
\item Penalize $O_g$, not raw use.
\item Give each reaction a functional-protection coefficient $c_r \in [0,1]$:
  high for shared membrane-conversion machinery, lower for
  identity-specific autocatalytic cycling.
\item Diffuse release pressure toward neighboring implementations via a
  motif-neighborhood kernel.
\end{enumerate}

In Hyperseed terms, this is \emph{fiber-selective eviction}: don't evict
fiber indiscriminately; evict only the excess --- fiber that exceeds its
value-justified share. Preserve shared functional fiber.

This maps directly to the attentional projection principle
(LTM 7aef7064): eviction must preserve epistemic integrity. In the
chemical context, ``epistemic integrity'' is ``membrane conversion
capacity'' --- the functional fiber that maintains the organism's
individuation. Evicting this fiber kills the organism, just as
evicting essential epistemic fiber kills the reasoner's coherence.

The species-agnostic character of the rule is essential: it doesn't
encode which implementation should win. It encodes the \emph{principle}
by which any implementation's dominance should be evaluated: value
contribution relative to ecological share. The principle is
fiber-general --- it applies to any system where implementations
compete for representational space.
\end{theorem}

%% =================================================================
\section{Identity vs.\ function as section vs.\ fiber}
\label{sec:identity}
%% =================================================================

\begin{definition}[The crucial distinction --- claims 18--20, 22]
\label{def:identity}
The article identifies the conceptual key: \emph{identity} (which
specific implementation occupies a niche) vs.\ \emph{function} (what
it does for the organism). Anti-precedence should target identity
lock-in, not functional repetition.

\begin{itemize}
\item \textbf{Identity} = the specific autocatalytic cycle (expendable,
  replaceable). A particular RAF set that happens to catalyze the right
  reactions.
\item \textbf{Function} = membrane conversion machinery (essential,
  must be preserved across identity changes). The shared infrastructure
  that maintains individuation.
\end{itemize}

In Hyperseed terms:
\begin{itemize}
\item \textbf{Identity = section.} A specific implementation is a
  section of the fiber bundle --- one particular assignment of
  implementations to functional roles. Sections are expendable:
  you can change which implementation occupies a role without
  changing the role itself.
\item \textbf{Function = fiber.} The functional role is the fiber ---
  the abstract structure that any section must instantiate. The fiber
  persists across section changes. Membrane conversion is a fiber
  property; which RAF set performs it is a section choice.
\end{itemize}

Anti-precedence targets sections (specific implementations), not fibers
(functional roles). This is the same section/fiber distinction that
appears throughout the Hyperseed ontology:
\begin{itemize}
\item In linguistics: specific linguistic forms (sections) vs.\
  cognitive functions they express (fibers).
\item In governance: specific office-holders (sections) vs.\
  governance functions (fibers).
\item In AI: specific model weights (sections) vs.\ capabilities
  they implement (fibers).
\end{itemize}
\end{definition}

%% =================================================================
\section{Conversion bridge as transition function}
\label{sec:bridge}
%% =================================================================

\begin{proposition}[Bridging old and new --- claim 23]
\label{prop:bridge}
The ``conversion bridge'' that enables transition from old implementation
to new is a transition function in the fiber bundle: it maps from the
old section to the new section while preserving the functional fiber.

The bridge must satisfy two constraints simultaneously:
\begin{enumerate}
\item \textbf{Functional preservation:} membrane conversion (functional
  fiber) must be maintained throughout the transition. The organism
  can't survive a gap in membrane production.
\item \textbf{Identity change:} the dominant autocatalytic cycle (section)
  must shift from old to new. The old implementation must yield share
  to the new one.
\end{enumerate}

In fiber bundle terms, this is a transition function $g_{12}$ that
maps from the old section $s_1$ to the new section $s_2$ while
preserving the fiber:
\[
  g_{12}: F|_{U_1} \to F|_{U_2} \qquad \text{with} \qquad
  \pi \circ g_{12} = \pi
\]

The bridge preserves the projection $\pi$ (functional output) while
changing the fiber assignment (which implementation produces that
output). This is why the conversion bridge works and the functional
bridge alone doesn't: the conversion bridge explicitly preserves
the functional fiber during the transition, while other bridges
might preserve structural similarity without guaranteeing functional
continuity.
\end{proposition}

%% =================================================================
\section{Biphasic memory failure as indiscriminate eviction}
\label{sec:failure}
%% =================================================================

\begin{proposition}[The wrong eviction rule --- claims 5--6, 27]
\label{prop:failure}
The broad biphasic memory rule failed because it penalized repetition
itself --- and repetition is how a chemical individual forms. Without
recurrence, there is no self-maintaining organization.

In Hyperseed terms, this is \emph{indiscriminate fiber eviction}:
treating all repetition as potentially problematic, which evicts
essential fiber (membrane conversion machinery) along with excess
fiber (over-canalized identity).

The lesson is general: any eviction/forgetting/anti-precedence rule
must distinguish between:
\begin{itemize}
\item \textbf{Essential repetition} (functional fiber): repetition
  that maintains the system's core capabilities. Evicting this
  destroys the system.
\item \textbf{Excess repetition} (over-canalized section): repetition
  that locks in a specific implementation beyond its justified share.
  Evicting this improves the system's adaptability.
\end{itemize}

The biphasic rule couldn't make this distinction because it operated
on raw repetition counts, not on value-relative shares. The protected
release rule can make the distinction because it separates the
\emph{function} question (``is this repetition functionally valuable?'')
from the \emph{identity} question (``is this implementation
over-represented?'').

This connects to the attentional projection failure mode: an attention
system that evicts based on recency or frequency alone (without
considering functional importance) will eventually evict essential
knowledge. The chemical experiment provides computational evidence
that this failure mode is real and that the fiber-selective alternative
works.
\end{proposition}

%% =================================================================
\section{Value-justified share as fiber-weighted allocation}
\label{sec:allocation}
%% =================================================================

\begin{definition}[Fiber budget management --- claims 8--9, 24--25]
\label{def:allocation}
Each pathway's share should be proportional to its fiber contribution
(value). Excess canalization is fiber over-representation: a pathway
taking up more fiber space than its contribution justifies.

In Hyperseed terms, the value-justified share is \emph{fiber-weighted
allocation}: the same principle as reputation-weighted governance in
OpenBGI, applied to chemical ecology:
\[
  S^*_g(t) = \frac{\exp(V_g(t) / \tau)}{\sum_{g'} \exp(V_{g'}(t) / \tau)}
\]
where $V_g(t)$ is the value (membrane-conversion payoff per unit cost)
of pathway family $g$ at time $t$.

The anti-canalization rule is a fiber-budget management principle:
the total fiber space (ecological capacity) is finite, and each
pathway should occupy a share proportional to its contribution.
Over-representation of low-value pathways crowds out high-value
alternatives.

The motif-neighborhood kernel (claim 26) ensures that release
pressure follows \emph{local transition functions}: it flows to
structurally similar alternatives, not random reassignment. This
preserves the topology of the implementation space --- nearby
implementations are more likely to be functionally compatible.
\end{definition}

%% =================================================================
\section{Cross-references}
%% =================================================================

\begin{remark}[Connections to other formalizations]
\begin{itemize}
\item \textbf{Attentional projection} (LTM 7aef7064): eviction
  integrity. The chemical experiment validates the principle that
  eviction must preserve functional fiber.
\item \textbf{OpenBGI} (2026-05-07): fiber-weighted governance.
  Value-justified share in chemistry parallels reputation-weighted
  governance: allocation proportional to contribution.
\item \textbf{Goals That Grow Back} (2026-07-20): anti-canalization.
  Goals that resist over-commitment to specific implementations
  --- the organizational analog of protected release.
\item \textbf{Seeding RSI} (2026-07-28): Hyperon architecture.
  The neural-symbolic architecture that could implement
  fiber-selective eviction in AI systems.
\item \textbf{The Orchard Bug} (2026-06-05): composition.
  The conversion bridge's success depends on maintaining
  the cocycle condition (functional continuity) during
  the transition.
\item \textbf{Grand Unified Physics} (2026-05-15): weakness
  principle. Value-relative anti-precedence is a form of
  weakness: the system favors implementations with least
  unnecessary dominance.
\end{itemize}
\end{remark}

\begin{conjecture}[Universal protected release]
Conjecture: the protected release principle (penalize excess
canalization, not repetition; protect shared function; release
toward neighbors) is universal across adaptive systems. In
chemistry, it governs evolutionary transitions between
autocatalytic sets. In AI, it should govern the replacement of
learned representations. In organizations, it should govern
leadership transitions. In ecosystems, it governs species
replacement.

If this conjecture holds, protected release is a fiber-bundle
invariant of adaptive systems: any system that maintains
individuation through functional repetition while remaining
adaptable through identity flexibility must implement some
form of fiber-selective eviction with functional protection.
\end{conjecture}

\end{document}
